% Generated by docs/report/make_result_tables.py -- do not edit by hand.
% Needs: booktabs (and a CJK-capable setup such as ctex).

\begin{table}[t]
\centering\small
\caption{模拟版图设计任务上现成 agent 的表现：论文报告值与本实验复现值。“完全通过”指全部检查项通过；“有效版图”指提交了可读取的版图文件；各项为通过的运行占全部运行的比例。}
\label{tab:task}
\vspace{3pt}
\footnotesize\setlength{\tabcolsep}{3.5pt}
\begin{tabular}{@{}lcccccccc@{}}
\toprule
模型 + agent 工具 & 总分 & 完全通过 & 有效版图 & DRC & LVS & 后仿真 & 工具调用 & Token (M) \\
\midrule
\multicolumn{9}{@{}l}{\textit{论文报告（40 题 $\times$ 3 次）}\textsuperscript{a}} \\
DeepSeek V4.1 Flash + Claude Code & 0.1974 & 8.3\% & -- & 41.7\% & 29.2\% & 17.5\% & 134 & 13.2 \\
GPT-6 Astra + Codex\textsuperscript{d} & 0.7483 & 70.8\% & -- & 99.2\% & 98.3\% & 82.5\% & 112 & 11.2 \\
Novara Agent System\textsuperscript{d} & 0.8154 & 72.5\% & -- & 97.5\% & 100.0\% & 75.0\% & -- & -- \\
\midrule
\multicolumn{9}{@{}l}{\textit{本实验（14 题 $\times$ 1 次）}} \\
DeepSeek V4.1 Flash + Claude Code\textsuperscript{b} & 0.0571 & 0.0\% & 42.9\% & 14.3\% & 7.1\% & 0.0\% & 288 & 26.3\textsuperscript{c} \\
\quad（通过次数/运行次数） & & (0/14) & (6/14) & (2/14) & (1/14) & (0/14) & & \\
\bottomrule
\end{tabular}

\vspace{2pt}
\begin{minipage}{\linewidth}\footnotesize
\textsuperscript{a} 取自论文表 1，为 40 道题各运行 3 次的平均；论文没有按配置报告版图提交率（只给出全部 13 种配置合计为 91.1\%），也未公布逐题结果。本实验使用其中 14 道题，各运行 1 次，题目范围与次数均与论文不同。\\
\textsuperscript{b} 14 次运行中有 8 次到结束时没有提交版图（记 0 分），其中 4 次是因为单次回复超过客户端的输出长度上限、会话被中止。对这 8 次的运行记录逐步分析的结果是：全程只做过器件级的试验，都没有生成过包含全部器件的版图。论文为自行部署模型，本实验通过 API 网关调用；与论文同配置结果的差距，原因尚未查清。\\
\textsuperscript{c} 每次运行的输入（含缓存命中）与输出 token 之和，单位百万。论文表注为“记录的输入与输出 token 累计”，未说明是否含缓存命中；不含缓存命中时本实验为 23.6。\\
\textsuperscript{d} 分别为论文中得分最高的通用 agent 和版图专用系统，列出以示该任务可达到的水平。\\ 说明：基准当前版本的后仿真只执行公开测试台中写明的工况（施密特触发器一题为题面 9 个工艺角与温度组合中的 1 个），“后仿真”和“完全通过”均按此口径。
\end{minipage}
\end{table}

\begin{table}[t]
\centering\small
\caption{对 agent 的两级改进在 14 道留出测试题上的效果（上），以及自动优化两轮的过程（下）。模型与 agent 工具始终不变，改动的只是其外部的辅助流程。}
\label{tab:improve}
\vspace{3pt}
\setlength{\tabcolsep}{5pt}
\begin{tabular}{@{}lcccccccc@{}}
\toprule
\multicolumn{9}{@{}l}{\textbf{(a) 留出测试集}（14 题，每种条件各运行 1 次）} \\
\midrule
条件 & 总分 & 较上一行 & 升/平/降 & 完全通过 & 有效版图 & DRC & LVS & 后仿真 \\
\midrule
原始 agent & 0.0571 & -- & -- & 0/14 & 6/14 & 2/14 & 1/14 & 0/14 \\
\quad + 人工编写的辅助流程\textsuperscript{a} & 0.1453 & +0.0882 & 8/5/1 & 1/14 & 14/14 & 10/14 & 1/14 & 1/14 \\
\quad\quad + 两轮自动优化\textsuperscript{b} & 0.1453 & +0.0000 & 0/14/0 & 1/14 & 14/14 & 11/14 & 1/14 & 1/14 \\
\bottomrule
\end{tabular}

\vspace{6pt}
\begin{tabular}{@{}clcccc@{}}
\toprule
\multicolumn{6}{@{}l}{\textbf{(b) 自动优化过程}（开发集 12 题，每题 2 次，共 24 次运行）} \\
\midrule
轮次 & 修改内容 & 尝试次数\textsuperscript{c} & 平均分变化 & 是否采纳 & 模型调用 \\
\midrule
1 & 修补环节的提示中加入各项检查的失败原因 & 6 & +0.0023 & 是 & 2792 \\
2 & 再次改写修补环节的提示（针对 LVS） & 3 & +0.0005 & 是 & 3032 \\
\bottomrule
\end{tabular}

\vspace{2pt}
\begin{minipage}{\linewidth}\footnotesize
\textsuperscript{a} 人工编写的辅助流程包括：规定在草稿目录中工作并可随时送检、交卷；第 140 分钟给出时间提醒；第 165 分钟停止主会话，由一个新会话做最多 15 分钟的修补。该行增益的 59\% 来自一道题（唯一完全通过的那道）。\\
\textsuperscript{b} 两轮被采纳的修改都只涉及最后的修补环节，因此该行是在上一行各次运行第 165 分钟的存档上重新执行修补环节得到的，此前的过程与上一行相同。\\
\textsuperscript{c} 自动提出修改方案的 agent 得到一个可评测的方案所用的尝试次数。\\
说明：开发集上“人工辅助流程”的平均分为 0.1640（24 次运行）。各条件在测试集上只运行了 1 次，没有重复实验，无法给出随机波动的范围。
\end{minipage}
\end{table}

\begin{table}[t]
\centering\small
\caption{14 道测试题在三种条件下的逐题结果。G：提交了有效版图；D：DRC 通过；L：LVS 通过；S：后仿真通过；“--”：没有提交版图。}
\label{tab:pertask}
\vspace{3pt}
\begin{tabular}{@{}lccccccc@{}}
\toprule
 & & \multicolumn{2}{c}{原始 agent} & \multicolumn{2}{c}{+ 人工辅助流程} & \multicolumn{2}{c}{+ 两轮自动优化} \\
\cmidrule(lr){3-4}\cmidrule(lr){5-6}\cmidrule(l){7-8}
电路 & 器件数 & 得分 & 通过项 & 得分 & 通过项 & 得分 & 通过项 \\
\midrule
八管施密特触发器 & 8 & 0.00 & -- & 0.73 & GDLS & 0.73 & GDLS \\
Beta-multiplier 电流基准 & 9 & 0.00 & -- & 0.10 & G & 0.10 & G \\
60 dB、200 MHz 两级米勒运放 & 12 & 0.10 & G & 0.10 & GD & 0.10 & G \\
两级米勒补偿 OTA & 12 & 0.00 & -- & 0.10 & G & 0.10 & GD \\
非稳压电荷泵 & 14 & 0.00 & -- & 0.10 & GD & 0.10 & GD \\
70 dB、20 MHz 电流偏置两级米勒运放 & 15 & 0.00 & -- & 0.10 & GD & 0.10 & GD \\
折叠级联 OTA & 20 & 0.00 & -- & 0.10 & GD & 0.10 & GD \\
可编程 ICC/IPTAT NMOS 电流镜 & 26 & 0.00 & -- & 0.10 & GD & 0.10 & GD \\
1 GHz 二分频触发器 & 32 & 0.10 & G & 0.10 & G & 0.10 & G \\
可使能 RC 振荡器 & 32 & 0.00 & -- & 0.10 & G & 0.10 & GD \\
电流饥饿环形 VCO & 34 & 0.10 & G & 0.10 & GD & 0.10 & GD \\
10 MHz 稳压电荷泵 & 35 & 0.10 & GD & 0.10 & GD & 0.10 & GD \\
130 dB增益增强折叠共源共栅OTA & 60 & 0.10 & G & 0.10 & GD & 0.10 & GD \\
900 MS/s 浮动电荷转移余量放大器 & 118 & 0.30 & GDL & 0.10 & GD & 0.10 & GD \\
\midrule
平均 & & 0.0571 &  & 0.1453 &  & 0.1453 &  \\
\bottomrule
\end{tabular}
\end{table}
